\setcounter{section}{61} % tema número N-1
\section{Sistemas moviles}

Nombre completo: Conceptos básicos de otros sistemas operativos: OSx, iOS, Android Z/OS. Sistemas operativos para dispositivos móviles

\localtableofcontents

\subsection{Contexto normativo}
    Hay dos normas relacionadas con dispositivos móviles:
    \begin{description}
        \item[RD 1112/2018] sobre accesibilidad a sitios web y dispositivos móviles del sector público.
        \item[App Factory] incluida en España Digital 2026 y en el plan de Digitalización de administraciones públicas
    \end{description}

\subsection{Mac OS}
    Desarrollado y comercializado por Apple en su propio hardware. Está desarrollado para plataformas x86 y ARM (y powerPC antes) sobre los procesadores de apple. Está basado en BSD y Unix por lo que resulta muy familiar en entornos unix e implementa POSIX

    OSx tiene diferentes estructuras que permiten trabajar a desarrolladores de diferentes comunidades: Cocoa permite el desarrollo de aplicaciones nativas y el uso de interfaz gráfica, Carbon simplifica la migración de aplicaciones para versionas anteriores a la 10.

\subsection{iOS}
    Desarrollado originalmente para iPhone, y mas tarde para los dispositivos móviles. está basado en OSx, por lo que hereda muchos elementos de este, está rediseñado para usarse con interfaces táctiles. El desarrollo de aplicaciones está limitado al entorno apple. Limitando al usuario la ejecución de aplicaciones no aprobadas.

\subsection{Android}
    Desarrollado por google y basado en Linux. Google realiza dos desarrollos paralelos, incluyendo la versión libre AOSP. Android basa sus aplicaciones en Java, desarrollando su propia máquina virtual, y mas adelante compilación Just In Time y Ahead of Time, traduciendo código java a bytecode.

\subsection{Z/OS}
    Es el sistema operativo creado por IBM para sus mainframes basado en 64 bits. Se centra en características para este tipo de entornos, como fiabilidad, disponibilidad, tolerancia a fallos, escalabilidad\dots